The land-surface model SURFEX significantly impacts the internal climatology of long-term climate simulations by improving the representation of land surface processes, which in turn affects temperature and precipitation extremes, urban heat islands, and overall climate projections.

Key Impacts of SURFEX on Climate Simulations

Improved Temperature and Precipitation Representation: SURFEX enhances the simulation of temperature and precipitation extremes by providing a more accurate representation of land-atmosphere interactions. This is crucial for reducing uncertainties in climate model projections, as different land surface models can lead to significant variations in the intensity and duration of extreme events (García‐García et al., 2020).


Urban Climate Modeling: SURFEX, when coupled with urban canopy models like TEB, significantly improves the simulation of urban heat islands (UHI) and urban climate processes. This is evident in studies over cities like Paris and Budapest, where SURFEX-TEB reduced biases in temperature and better captured the UHI effect compared to other models (Nogueira et al., 2022; Zsebeházi & Szépszó, 2020).


Vegetation and Land Cover Representation: SURFEX's use of updated land cover datasets and interactive vegetation models helps in accurately simulating land surface temperatures and vegetation cover, which are critical for surface-atmosphere energy exchanges. This leads to better climate simulations, as seen in the Iberian Peninsula case study (Nogueira et al., 2020).


Regional Climate Simulations: The model's ability to be reinitialized daily or run continuously allows for flexibility in simulating regional climates. This approach has shown improvements in temperature simulations and energy flux partitioning, particularly during summer when land-atmosphere interactions are strong (Berckmans et al., 2016).


Adaptability and Versatility: SURFEX can be used in various configurations, including offline and coupled modes, making it suitable for a wide range of applications from mesoscale models to global climate simulations. This adaptability helps in mutualizing research efforts and improving forecast performance (Hamdi et al., 2013; Sørland et al., 2021).

Conclusion

SURFEX significantly enhances the internal climatology of long-term climate simulations by improving the accuracy of temperature and precipitation extremes, urban climate modeling, and land-atmosphere interactions. Its adaptability and detailed representation of land surface processes make it a valuable tool for climate research and projections.

These papers were sourced and synthesized using Consensus, an AI-powered search engine for research. Try it at https://consensus.app

References

Nogueira, M., Albergel, C., Boussetta, S., Johannsen, F., Trigo, I., Ermida, S., Martins, J., & Dutra, E. (2020). Role of vegetation in representing land surface temperature in the CHTESSEL (CY45R1) and SURFEX-ISBA (v8.1) land surface models: a case study over Iberia. Geoscientific Model Development. https://doi.org/10.5194/gmd-2020-49-supplement

Berckmans, J., Giot, O., Troch, R., Hamdi, R., Ceulemans, R., & Termonia, P. (2016). Reinitialised versus continuous regional climate simulations using ALARO-0 coupled to the land surface model SURFEXv5. Geoscientific Model Development, 10, 223-238. https://doi.org/10.5194/GMD-10-223-2017

Nogueira, M., Hurduc, A., Ermida, S., Lima, D., Soares, P., Johannsen, F., & Dutra, E. (2022). Assessment of the Paris urban heat island in ERA5 and offline SURFEX-TEB (v8.1) simulations using METEOSAT land surface temperature product. Geoscientific Model Development. https://doi.org/10.5194/gmd-2021-431

Zsebeházi, G., & Szépszó, G. (2020). Modeling the urban climate of Budapest using the SURFEX land surface model driven by the ALADIN-Climate regional climate model results. **, 124, 191-207. https://doi.org/10.28974/idojaras.2020.2.3

García‐García, A., Cuesta-Valero, F., Beltrami, H., González-Rouco, F., García-Bustamante, E., & Finnis, J. (2020). Land surface model influence on the simulated climatologies of temperature and precipitation extremes in the WRF v3.9 model over North America. Geoscientific Model Development, 13, 5345-5366. https://doi.org/10.5194/gmd-13-5345-2020

Hamdi, R., Degrauwe, D., Duerinckx, A., Cedilnik, J., Costa, V., Dalkilic, T., Essaouini, K., Jerczynki, M., Kocaman, F., Kullmann, L., Mahfouf, J., Meier, F., Sassi, M., Schneider, S., Váňa, F., & Termonia, P. (2013). Evaluating the performance of SURFEXv5 as a new land surface scheme for the ALADINcy36 and ALARO-0 models. Geoscientific Model Development, 7, 23-39. https://doi.org/10.5194/GMD-7-23-2014

Sørland, S., Brogli, R., Pothapakula, P., Russo, E., Van De Walle, J., Ahrens, B., Anders, I., Bucchignani, E., Davin, E., Demory, M., Dosio, A., Feldmann, H., Früh, B., Geyer, B., Keuler, K., Lee, D., Li, D., Van Lipzig, N., Min, S., Panitz, H., Rockel, B., Schär, C., Steger, C., & Thiery, W. (2021). COSMO-CLM regional climate simulations in the Coordinated Regional Climate Downscaling Experiment (CORDEX) framework: a review. Geoscientific Model Development. https://doi.org/10.5194/gmd-14-5125-2021